\section{Related Work}
\label{sec:related-works}

\paragraph{Layer pruning and skipping.} A large body of work removes Transformer
blocks outright, differing mainly in the importance score used to choose them:
angular distance between block boundaries \citep{gromov2024unreasonable}, block
influence \citep{men2024shortgpt}, structured dropout that trains the model to be
prunable \citep{fan2020reducing}, and gradient- or activation-based structural
pruning \citep{ma2023llmpruner, frantar2023sparsegpt, sun2024simple}. Dynamic
variants route tokens around blocks at inference time
\citep{raposo2024mixture, elhoushi2024layerskip, schuster2022confident}. These
methods answer \emph{which} block to drop. Once the choice is made, essentially
all of them apply the same operator, set the residual update to zero, and it
is that operator, not the selection rule, that \mname replaces. The two are
orthogonal: \mname consumes whatever layer set a selection method produces.

\paragraph{Recovering what pruning removes.} Post-pruning repair typically fits a
map at the boundary of the removed region: a linear transformation between hidden
states at different depths \citep{din2023jump}, or a learned lens that reads
predictions out of intermediate states \citep{belrose2023eliciting}. Such a map
costs $O(d^2)$ parameters and $O(Td^2)$ work, the same order as the block it
replaces, which erodes the saving that motivated pruning. \mname instead
predicts from the \emph{preceding update trajectory} rather than from the
boundary activation, which is what allows the predictor to collapse to one or two
scalars and $O(Td)$ work.

\paragraph{Depth as a trajectory.} That residual streams evolve smoothly is well
established: residual networks behave like ensembles of shallow paths
\citep{veit2016residual}, and residual connections induce iterative refinement of
a single representation across depth \citep{jastrzebski2018residual}. This
literature is largely analytical. We take it as an engineering resource: if
updates lie on a smooth path, the next one is forecastable from the last, and a
skipped block need not be a hole in the computation.

To our knowledge, \mname is the first layer-skipping method to model the missing
residual update autoregressively from the preceding depth trajectory using a
closed-form, per-channel (diagonal) predictor, and, in the same breath, the first to
show that doing so recovers likelihood without recovering function.
